
\sgaStableTarget{sga6.e10.sec.7.5}\subsection*{7.5. ``Conservation of number''.}
Suppose now that @@M0@@ is proper, so that @@M1@@ and @@M2@@ are proper over @@M3@@ and @@M4@@ respectively, and one has the homomorphisms of 2.3:
@@M5@@
Then one has the relation
\sgaStableTarget{sga6.e10.eq.7-5-1}@@M6@@
(the ``conservation-of-number theorem''); in other words, the diagram
\sgaStableTarget{sga6.e10.eq.7-5-2}@@M7@@
commutes. To prove (7.5.1), one is again reduced to the case where @@M8@@, with @@M9@@ an integral closed subscheme of @@M10@@, and, with the notation of the proof of 7.4, one is reduced to proving that
@@M11@@
which is well known (EGA III 7.9.4, or Exposé III 5.7).

\sgaStableTarget{sga6.e10.sec.7.6}\subsection*{7.6.}
It is immediate that
\sgaStableTarget{sga6.e10.eq.7-6-1}@@M12@@
which, taking account of the @@M14@@-linearity of @@M13@@, implies that the following diagram is commutative:
\sgaStableTarget{sga6.e10.eq.7-6-2}@@M15@@
